\chapter{Tools and Technology Used}

\section{Software Technologies}
OmniFAQ AI is built using a modern, decoupled web technology stack that delivers exceptional responsiveness, high security, and minimal external dependency overhead.

\subsection{Frontend}
\begin{itemize}
    \item \textbf{HTML5:} Structures the Single Page Application, defining semantic containers, authentication drawers, search interfaces, and question card grids.
    \item \textbf{CSS3 (Glassmorphism):} Provides a sleek, dark-slate teal design system (\texttt{\#112F3B}) with vibrant cyan accents (\texttt{\#0891B2}, \texttt{\#06B6D4}), backdrop blur filters, custom scrollbars, and responsive layouts. The entire core stylesheet is inlined to ensure zero render-blocking layout shift.
    \item \textbf{JavaScript / jQuery 3.7.1:} Orchestrates client-side view routing, search input debouncing, asynchronous REST API communication via AJAX, accordion-style answer reveals, and active recall practice scoring.
    \item \textbf{Lucide Icons:} Supplies lightweight, high-contrast SVG vector icons across navigation headers, difficulty badges, and input controls.
\end{itemize}

\subsection{Backend}
\begin{itemize}
    \item \textbf{PHP 8.2:} Functions as the core server-side runtime, providing a centralized RESTful router (\texttt{api.php}) with CORS support, JSON serialization, and role-based access validation.
    \item \textbf{Dual-Driver Database Engine (MySQL \& SQLite):} Connects to cloud MySQL/MariaDB instances when environment variables are supplied, with automatic zero-configuration failover to an embedded SQLite database (\texttt{database.sqlite}).
    \item \textbf{Google Gemini API (\texttt{gemini-2.5-flash}):} Delivers advanced generative intelligence for synthesizing novel technical interview questions under strict JSON schema constraints.
    \item \textbf{Docker \& Apache 2.4:} Encapsulates the application into a portable, production-ready container deployed seamlessly on cloud PaaS platforms like Render.
\end{itemize}

\section{System Architecture}
The system follows a three-tier client-server architecture:
\begin{enumerate}
    \item \textbf{Presentation Tier:} The browser client handles user interactions, manages session tokens in \texttt{localStorage}, and emits asynchronous HTTP requests.
    \item \textbf{Application Tier:} The Apache web server and PHP router parse requests, execute authentication routines, perform weighted relevance ranking, and handle outbound cURL requests to Google's API.
    \item \textbf{Data Tier:} The relational database manages users, FAQs, and search telemetry, while the Gemini API serves as an external generative knowledge engine.
\end{enumerate}

\section{Software Model Used}
The project was developed following the \textbf{Iterative Development Model}. An initial prototype supported basic question querying and local SQLite storage. Subsequent iterations incrementally integrated real-time search debouncing, Google Gemini API synthesis, automated database persistence, active recall practice evaluation, and cloud Docker containerization.

\section{Database Design}
The relational schema comprises three primary tables: \texttt{users}, \texttt{faqs}, and \texttt{search\_analytics}.

\subsection{Important Fields}
\begin{itemize}
    \item \textbf{\texttt{users} Table:}
    \begin{itemize}
        \item \texttt{id}: Primary Key, Auto-increment.
        \item \texttt{full\_name}: String, candidate display name.
        \item \texttt{email}: String, Unique identifier used for login.
        \item \texttt{password}: String, Bcrypt salted hash.
        \item \texttt{role}: String (\texttt{user} or \texttt{admin}).
    \end{itemize}
    \item \textbf{\texttt{faqs} Table:}
    \begin{itemize}
        \item \texttt{id}: Primary Key, Auto-increment.
        \item \texttt{question}: Text, the technical interview question.
        \item \texttt{answer}: Text, comprehensive model answer.
        \item \texttt{category}: String (JavaScript, React, Node.js, SQL, System Design).
        \item \texttt{difficulty}: String (Easy, Medium, Hard).
        \item \texttt{tags}: String, comma-separated keywords for relevance matching.
        \item \texttt{views}: Integer, tracks question popularity.
    \end{itemize}
    \item \textbf{\texttt{search\_analytics} Table:}
    \begin{itemize}
        \item \texttt{id}: Primary Key, Auto-increment.
        \item \texttt{query\_text}: String, search query entered by the user.
        \item \texttt{matched\_results\_count}: Integer, count of matching FAQs returned.
        \item \texttt{is\_unresolved}: Integer (1 if 0 matches found; 0 otherwise).
        \item \texttt{searched\_at}: Timestamp of search execution.
    \end{itemize}
\end{itemize}

\section{Product User Interfaces}
The primary interface is a search-first Single Page Application consisting of:
\begin{itemize}
    \item A prominent debounced search bar with one-click suggestion tags (JavaScript, React, Node.js, SQL, System Design).
    \item An interactive question card grid displaying domain badges, difficulty tiers, view counts, and collapsible model answers.
    \item An active recall practice text area with automated keyword coverage scoring.
    \item A secured administrative analytics modal displaying platform KPIs and top unresolved student queries.
\end{itemize}

\section{Design and Implementation Constraints}
\begin{itemize}
    \item \textbf{Dual Database Portability:} The system must run flawlessly in both zero-config local classrooms (using SQLite) and scalable cloud clusters (using MySQL) without altering business logic code.
    \item \textbf{API Key Security:} Google Gemini API keys must never be exposed to the client-side JavaScript bundle; all AI transactions are securely handled server-side through PHP cURL routines and environment variables.
    \item \textbf{Dynamic Port Allocation:} The Docker container must dynamically adapt its Apache listening port at runtime to comply with cloud PaaS hosting environments.
\end{itemize}
